\section{Mathematical typography in \LaTeX}

	\subsection{Superscripts and subscripts}
		Superscripts and subscripts are commonly used to represent exponents (such as $3^2$, $x^{-3}$) and indexed quantities (such as $m_1$, $K_0$) respectively. In \LaTeX, superscripts are typeset using the \emph{caret} symbol (\lstinline|^|) while subscripts are typeset using the \emph{underscore} symbol (\lstinline|_|).
		
		If a superscript or subscript contains more than a single character, the entire exponent or index must be enclosed in braces \lstinline|{ }|.
		
		For instance, the following code contains expressions with exponents and indices
\lstinputlisting{codes/snippet-sup-sub}
		which produces the output
		\lstoutput{codes/snippet-sup-sub}
	
	\subsection{Fractions and roots}
	
	\subsection{Pre-defined functions and operators}
	
	\subsection{Summation, product and integral operators}
	
	\subsection{Greek letters}
	
	\subsection{Mathematical fonts}